%---------------------------------------------------------------
\chapter{Návrh řešení}
\label{ch:navrh}
%---------------------------------------------------------------

\section{Architektura systému}
\label{sec:architektura}
% TODO: Vue.js 3 + FastAPI + Celery + Redis + PostgreSQL
% - Diagram architektury (figure)
% - Upload flow: video + úroveň + pohled + styl → WebSocket progress → výsledky
% - Asynchronní zpracování (Celery workers)

\section{Výběr modelu pro odhad pózy}
\label{sec:vyber_modelu}
% TODO: Porovnání MediaPipe vs RTMPose vs ViTPose
% - Důvody výběru (přesnost, rychlost, fine-tuning schopnosti)
% - Tabulka: model × AP × FPS × parametry

\section{Ragdoll constraints}
\label{sec:ragdoll}
% TODO: Post-processing vrstva po pose estimation
% - Limity rotace kloubů (loket 0-150°, koleno 0-170°, atd.)
% - Filtrování fyzikálně nemožných póz (artefakty z okluzí, reflexů)
% - Temporální vyhlazení (výkyvy mezi framy)
% - Detekce záměny levá/pravá strana
% - Důležitost pro podvodní záběry

\section{Biomechanické metriky}
\label{sec:metriky}
% TODO: Úhly kloubů, rotace trupu, frekvence/délka záběru, fáze záběru
% - Tabulka: metrika × úroveň × optimální rozsah × práh chyby
% - 3 úrovně: začátečník / mírně pokročilý / expert

\section{Referenční systém}
\label{sec:referencni_system}
% TODO: 3 úrovně × pohledy × styly
% - Generování referenčních vzorů ze SwimXYZ
% - DTW porovnání (vysvětlení algoritmu)

\section{Porovnávací metoda}
\label{sec:porovnavaci_metoda}
% TODO: DTW + pravidlový systém + skóre 0-100
% - Detekce cyklů záběru
% - Výpočet odchylek od reference
% - Mapování odchylek na konkrétní chyby
% - Generování feedbacku v češtině

\section{Diskuse alternativních přístupů}
\label{sec:diskuse_alternativ}
% TODO: RAG/LLM přístup — proč jsme se rozhodli PROTI
% - Comendant 2024 (RAG coaching for swimming)
% - Talking Tennis (biomechanics → LLM)
% - BoxingPro (IoT → LLM)
% Důvody PROTI:
% - Závislost na externích API (single point of failure)
% - Nedeterministické výstupy (stejný vstup → různý feedback)
% - Latence a náklady na inference
% - Nemožnost plně fungovat offline
% - Obtížná reprodukovatelnost experimentů
% Trade-offs: vyšší flexibilita LLM vs. konzistence a kontrola pravidlového systému
